% ============================================================================
%  SECCIÓN 3.27: METODOLOGÍA DE DESARROLLO - SCRUM
% ============================================================================

\section{Metodología de desarrollo: Scrum}

Para el desarrollo del proyecto se adoptó la metodología ágil \textbf{Scrum},
que permite organizar el trabajo en iteraciones cortas (\textit{sprints}) con
entregas incrementales de valor \citep{schwaber2020, beck2001}. Esta elección
responde a la necesidad de obtener un MVP funcional en un plazo reducido de
tres semanas, aprovechando la flexibilidad de Scrum para adaptar los
requerimientos durante el desarrollo.

\subsection{Equipo Scrum}

El equipo se organizó con los siguientes roles:

\begin{table}[H]
  \centering
  \caption{Roles del equipo Scrum}
  \contenidoConFuente{%
    \begin{tablaAPA}
      \begin{tabular}{@{}p{4.5cm}m{10.5cm}@{}}
        \toprule
        \textbf{Rol} & \textbf{Responsabilidades} \\
        \midrule
        Product Owner & Define las prioridades del producto, valida las
        historias de usuario y aprueba las entregas. \\
        \addlinespace
        Scrum Master & Facilita las ceremonias, elimina bloqueos del equipo y
        asegura el cumplimiento del proceso. \\
        \addlinespace
        Development Team & Desarrolla, prueba y entrega los incrementos del
        producto en cada sprint. \\
        \bottomrule
      \end{tabular}
    \end{tablaAPA}%
  }{Elaboración propia con base a \citet{schwaber2020}.}
\end{table}

\subsection{Organización del equipo}

\begin{table}[H]
  \centering
  \caption{Integrantes y roles en el proyecto}
  \contenidoConFuente{%
    \begin{tablaAPA}
      \begin{tabular}{@{}p{4.5cm}m{3.0cm}m{8.0cm}@{}}
        \toprule
        \textbf{Integrante} & \textbf{Rol Scrum} & \textbf{Responsabilidad
        técnica} \\
        \midrule
        Joaquin Felipez Rojas & Scrum Master & Coordinación, revisión de
        código, documentación final. \\
        \addlinespace
        Nicolas Reguerin Meneses & Development Team & Backend: Supabase,
        base de datos, autenticación. \\
        \addlinespace
        David Cruz Huanca & Development Team & Navegación, mapa, filtros,
        búsqueda, manejo de estado. \\
        \addlinespace
        Angel Rojas Hinojosa & Development Team & UI/UX: pantallas,
        componentes reutilizables, testing. \\
        \bottomrule
      \end{tabular}
    \end{tablaAPA}%
  }{Elaboración propia.}
\end{table}

\subsection{Duración y estructura de sprints}

El proyecto se planificó en \textbf{3 sprints} de una semana cada uno, con un
total de 12 días hábiles de desarrollo:

\begin{table}[H]
  \centering
  \caption{Estructura de sprints del proyecto}
  \contenidoConFuente{%
    \begin{tablaAPA}
      \begin{tabular}{@{}p{2.5cm}m{2.5cm}m{5.5cm}m{4.0cm}@{}}
        \toprule
        \textbf{Sprint} & \textbf{Duración} & \textbf{Objetivo} &
        \textbf{Entregable} \\
        \midrule
        Sprint 1 & Días 1--3 & Configurar Supabase, esquema de base de datos,
        autenticación y datos semilla & App conectada a Supabase con auth
        funcional \\
        \addlinespace
        Sprint 2 & Días 4--8 & Implementar todas las pantallas del MVP con
        datos reales & App navegable: mapa, catálogo, detalle, favoritos \\
        \addlinespace
        Sprint 3 & Días 9--12 & Tema visual, corrección de bugs, pruebas y
        documentación & MVP pulido y listo para entregar \\
        \bottomrule
      \end{tabular}
    \end{tablaAPA}%
  }{Elaboración propia.}
\end{table}

\subsection{Ceremonias Scrum}

Se implementaron las ceremonias estándar de Scrum adaptadas al contexto
académico:

\begin{table}[H]
  \centering
  \caption{Ceremonias Scrum implementadas}
  \contenidoConFuente{%
    \begin{tablaAPA}
      \begin{tabular}{@{}p{3.5cm}m{2.5cm}m{2.5cm}m{6.0cm}@{}}
        \toprule
        \textbf{Ceremonia} & \textbf{Frecuencia} & \textbf{Duración} &
        \textbf{Descripción} \\
        \midrule
        Sprint Planning & Inicio de sprint & 30--45 min & Se definen las
        historias de usuario del sprint y se asignan tareas. \\
        \addlinespace
        Daily Standup & Diaria (async) & 5 min & Cada integrante reporta:
        qué hizo, qué hará y si tiene bloqueos. \\
        \addlinespace
        Sprint Review & Fin de sprint & 30 min & Demostración del
        incremento terminado al equipo y al docente. \\
        \addlinespace
        Retrospectiva & Fin de sprint & 15 min & Identificación de mejoras
        para el siguiente sprint. \\
        \bottomrule
      \end{tabular}
    \end{tablaAPA}%
  }{Elaboración propia con base a \citet{schwaber2020}.}
\end{table}

\subsection{Backlog del producto}

El Product Backlog se gestionó mediante \textbf{GitHub Projects} con un tablero
kanban. Las tareas se identificaron con códigos del tipo \texttt{T-XXX} y se
organizaron en las siguientes categorías:

\begin{itemize}
  \item \textbf{Infraestructura:} Configuración de entorno, dependencias,
    base de datos.
  \item \textbf{Funcionalidad:} Pantallas, navegación, autenticación,
    mapa, catálogo, favoritos.
  \item \textbf{Integración:} Conexión frontend--backend, hooks de datos,
    estados de carga.
  \item \textbf{Pulido:} Tema visual, corrección de bugs, pruebas,
    documentación.
\end{itemize}

\subsection{Burndown por sprint}

El progreso de cada sprint se midió con el tablero kanban de GitHub Projects,
donde las tareas avanzan por las columnas: \textit{To Do} $\rightarrow$
\textit{In Progress} $\rightarrow$ \textit{In Review} $\rightarrow$
\textit{Done}. A continuación se presenta un resumen de las tareas completadas
por sprint:

\begin{table}[H]
  \centering
  \caption{Tareas completadas por sprint}
  \contenidoConFuente{%
    \begin{tablaAPA}
      \begin{tabular}{@{}p{3.0cm}m{2.5cm}m{2.5cm}m{8.0cm}@{}}
        \toprule
        \textbf{Sprint} & \textbf{Creadas} & \textbf{Completadas} &
        \textbf{Principales entregas} \\
        \midrule
        Sprint 1 & 15 & 15 & Proyecto Supabase, esquema SQL, auth,
        pantallas login/registro, splash \\
        \addlinespace
        Sprint 2 & 24 & 24 & Hooks de datos, mapa, catálogo, detalle,
        favoritos, filtros, búsqueda \\
        \addlinespace
        Sprint 3 & 17 & 17 & Tema visual, componentes base, pruebas,
        documentación, entrega \\
        \bottomrule
      \end{tabular}
    \end{tablaAPA}%
  }{Elaboración propia en base al tablero de GitHub Projects.}
\end{table}

\subsection{Definición de Hecho (DoD)}

Una tarea se consideraba ``hecha'' cuando cumplía con los siguientes criterios:

\begin{enumerate}
  \item El código está completo y funcional en el dispositivo de prueba.
  \item Se realizó \textit{code review} por al menos un integrante del equipo.
  \item No genera errores en la consola del desarrollador.
  \item Se verificó en al menos un dispositivo Android físico.
  \item El código fue mergeado a la rama \texttt{develop} sin conflictos.
  \item Se documentó el cambio en el mensaje de commit con el formato
    \texttt{feat(scope): descripción}.
\end{enumerate}

\subsection{Adaptaciones al contexto académico}

Dado el plazo reducido de tres semanas, se realizaron las siguientes
adaptaciones al proceso estándar de Scrum:

\begin{itemize}
  \item Los sprints se comprimieron a una semana en lugar de las dos habituales.
  \item Las reuniones diarias se realizaron de forma asíncrona por chat,
    evitando reuniones innecesarias.
  \item Se utilizó herramientas de inteligencia artificial para generar el
    código base, reduciendo los tiempos de desarrollo.
  \item El Product Owner y el Scrum Master fueron el mismo integrante,
    simplificando la toma de decisiones.
\end{itemize}
